\section{Introduction}
\label{sec:intro}

Coarse-grained reconfigurable arrays (CGRAs) occupy the middle ground between
the flexibility of software and the efficiency of fixed-function
hardware~\cite{cgra-survey}. Rather than reconfiguring at the bit level, as an
FPGA does, a CGRA exposes a grid of word-level processing elements (PEs) whose
operations and interconnect are set by a compact configuration, allowing the
same silicon to be re-purposed at run time for different kernels.

This project implements such a fabric on an FPGA and uses it as a
\emph{runtime-reconfigurable accelerator} attached to a host computer over a
serial link. The design goals are threefold: (i)~a small but expressive fabric
whose behaviour is fully controlled, cycle by cycle, from the host; (ii)~a
robust host/device protocol that can be exercised and verified without the
board; and (iii)~a UNIX-style software stack that lets a user describe and run
accelerated kernels declaratively.

A single architectural decision runs through the whole system and gives the
report its narrative spine: each PE owns exactly \emph{one} register and
external operands enter the mesh \emph{only} at two edges. We trace the
consequences of that decision from the datapath up to the dataflows a user can
express and down to the bytes on the wire.

Concretely, this report contributes:
\begin{itemize}
  \item a fully host-controlled $4\times4$ CGRA fabric whose every step is
        gated by a compact, checksum-protected UART protocol
        (Sections~\ref{sec:arch}--\ref{sec:proto});
  \item a mapping of element-wise, reduction and weight-stationary systolic
        dataflows onto the array, with an explicit argument for \emph{why} the
        single-register/edge-injection constraint admits some and excludes
        others (Section~\ref{sec:dataflow});
  \item a layered host stack---a thin C library, a byte-accurate device model
        that serves as a golden reference, and a declarative command-line
        front end---that runs and is verified end to end without an FPGA
        (Sections~\ref{sec:sw},~\ref{sec:results});
  \item a wire-level cost model that shows the accelerator's observed speed is
        bounded by serial round trips rather than by the array
        (Section~\ref{sec:cost}).
\end{itemize}

The remainder of this report is organised around the logical model of the
system. Section~\ref{sec:bg} places the fabric in the CGRA and systolic
literature. Section~\ref{sec:arch} describes the processing element and the mesh
interconnect. Section~\ref{sec:proto} presents the control finite-state machine
and the UART protocol. Section~\ref{sec:dataflow} discusses how element-wise,
reduction and systolic dataflows map onto the array, including the
architectural constraint that shapes those mappings. Section~\ref{sec:sw}
outlines the host software stack, Section~\ref{sec:cost} models its wire-level
cost, and Section~\ref{sec:results} summarises the verification methodology.
